% Submission to RV 2026 (short paper, up to 8 pages excl. references).
% LNCS class.
\documentclass[runningheads]{llncs}

% Packages and macros shared across all sections.

\usepackage[T1]{fontenc}
\usepackage{graphicx}
\usepackage{amsmath,amssymb}
\usepackage{mathpartir}
\usepackage{xcolor}
\usepackage{listings}
\usepackage{hyperref}
\renewcommand\UrlFont{\color{blue}\rmfamily}
\urlstyle{rm}

\lstdefinelanguage{move}{
  morekeywords={module,struct,enum,fun,public,let,mut,if,else,match,return,
                use,as,vector,move,copy,native,has,drop,key,store,copy},
  morecomment=[l]{//},
  sensitive=true,
}
\lstset{
  language=move,
  basicstyle=\ttfamily\small,
  keywordstyle=\bfseries,
  commentstyle=\itshape\color{gray!60!black},
  showstringspaces=false,
  columns=fullflexible,
  numbers=none,
  xleftmargin=1em,
}

\newcommand{\apt}{\textsc{apt}}
\newcommand{\local}[1]{\mathrm{loc}(#1)}
\newcommand{\glob}[1]{\mathrm{glb}(#1)}
\newcommand{\param}[1]{\mathrm{par}(#1)}
\newcommand{\node}[2]{\langle #1, #2 \rangle}
\newcommand{\Live}{\checkmark}
\newcommand{\Pois}{\dagger}
\newcommand{\refn}[1]{\rho_{#1}}
\newcommand{\mut}{\mathsf{mut}}
\newcommand{\imm}{\mathsf{imm}}
\newcommand{\Anc}{\mathcal{A}}
\newcommand{\Desc}{\mathcal{D}}


\begin{document}

\title{Runtime Reference Safety via Access Path Invalidation}
\titlerunning{Runtime Reference Safety via Access Path Invalidation}

\author{Vineeth Kashyap\thanks{Corresponding author.}\inst{1} \and
        Wolfgang Grieskamp\inst{1} \and
        Teng Zhang\inst{1}}
\authorrunning{V. Kashyap et al.}
\institute{Aptos Labs, Palo Alto, USA\\
\email{vineeth@aptoslabs.com}}

\maketitle

\begin{abstract}
Move's bytecode verifier statically enforces a small set of reference
safety rules: references should not dangle, reads through an
immutable reference should not observe a value that has been changed
through another handle, and mutable references at function call
boundaries should not alias. In several contexts the static check is
unavailable, untrusted, or unnecessarily conservative: bytecode that
bypasses the verifier (in fuzzing or mutated-bytecode harnesses,
behind dynamic dispatch), operations whose body lives outside the
bytecode language (native code), and programs the verifier rejects
even though the executions they actually take would not violate any
rule. We describe a runtime discipline that enforces Move's
reference-safety rules during execution. The core data structure is
the \emph{Access Path Tree} (\apt): a forest, built lazily during
execution, whose nodes record the access paths along which each live
reference was derived. A single invalidation rule poisons references
whose paths have been disturbed by destructive events; the rule
covers destructive writes, moves, freezes, variant mutations, calls,
returns, and opaque effects uniformly, and call boundaries
additionally lock subtrees so that mutable arguments do not overlap.
The discipline is path-sensitive---it runs against the actual
execution trace, rather than against a join over branches---and is
designed to accept every trace the static check would.  We describe
the model, sketch its small-step rules, and discuss its relation to
per-location dynamic alias models such as Stacked Borrows and Tree
Borrows.
\keywords{Runtime verification \and Reference safety \and Aliasing
  \and Access paths \and Move.}
\end{abstract}

\section{Introduction}\label{sec:intro}

Move is a bytecode language with first-class references over local
and globally-stored values. Each function's bytecode passes through
a static reference-safety verifier, which approves the function only
if, on every reachable execution path, references stay in scope, do
not violate type safety through reads or writes, do not overlap as
mutable arguments at call boundaries, and do not allow an immutable
view to observe a value rewritten through another handle. The
verifier rejects bytecode that violates these rules.

We are interested in enforcing the same rules \emph{at execution
time}. There are several reasons one would want this:
\begin{itemize}
\item Fuzzing and mutated-bytecode harnesses deliberately exercise
inputs the verifier would reject, to find bugs in everything
downstream.
\item Dynamic dispatch and native code escape static analysis: the
verifier cannot see what happens through a dispatched function
pointer or inside a Rust-implemented native, but the reference
rules still need to hold across those boundaries.
\item The static check has to be sound, so it sometimes rejects
programs that would not actually violate any rule on the executions
they take. A runtime check sees the executions and can be more
precise.
\end{itemize}

The following Move fragment illustrates the kind of subtle case at
stake:
\begin{lstlisting}
enum E { V1(T1), V2(T2) }
let v: E = ...;
let inner: &mut T1 = match &mut v { V1(p) => p };
... // some path mutates v to be V2
*inner    // unsafe
\end{lstlisting}
The reference \texttt{inner} is in scope. Its lifetime has not ended.
Its declared type is unchanged. But the program has switched
\texttt{v} from \texttt{V1} to \texttt{V2}, so the substructure that
\texttt{inner} pointed into is no longer there. Dereferencing
\texttt{inner} now reads memory that has been recycled into a value
of an unrelated type.

A sound static analysis catches this only conservatively: it does
not know which branch will run, so it must assume the worst. By the
time \texttt{*inner} executes, the runtime has already observed which
branch ran, which mutation occurred, and which reference was derived
from which. With the right slice of that information recorded, we
can answer ``is this reference still safe to use?'' precisely. The
challenge is to find a representation that is small enough to be
cheap to maintain, and uniform enough to cover the various events
that can affect a reference.

\paragraph{This paper.}
We organize live references by the path of borrows that produced
them, using a hierarchical structure we call an \emph{Access Path
Tree} (\apt). Each live reference is attached to a node of some APT,
with edges labelled by field projections, variant projections, and
a single label for vector indices. A single rule decides when a
reference becomes unusable: when an event destructively changes the
substructure reachable through some node, every reference attached
to a related node is marked \emph{poisoned} for the rest of the
execution. Destructive writes, moves and overwrites, freezes,
variant mutations, calls, returns, and opaque effects all fit this
pattern; what varies between them is which set of related nodes is
affected and which mutability filter applies. Calls additionally
lock subtrees so that two mutable arguments cannot overlap.

\paragraph{Contributions.}
\begin{itemize}
\item A runtime discipline for Move's reference-safety rules,
  organized around access path trees and one invalidation rule
  (Sects.~\ref{sec:apt}, \ref{sec:semantics}).
\item A small-step presentation in which the various
  reference-affecting events are instances of the same rule
  (Sect.~\ref{sec:semantics}).
\item A discussion of what the runtime trace buys over a join-based
  static check (Sect.~\ref{sec:precision}) and how the discipline
  relates to per-location dynamic alias models such as Stacked
  Borrows and Tree Borrows~\cite{stackedborrows,treeborrows}
  (Sect.~\ref{sec:related}).
\end{itemize}

\section{Motivating Examples}\label{sec:examples}

The enum example in Sect.~\ref{sec:intro} is the headline case:
mutation to a parent value invalidates a reference still in scope
below it. Two further examples make the design constraints concrete.

\paragraph{Order matters.}
Here are two fragments that differ only in the order of their last
two statements:
\begin{lstlisting}
let x = 0; let r2 = &mut x; let r1 = &mut x;       // (a)
*r2 = 5; *r1
let x = 0; let r2 = &mut x; let r1 = &mut x;       // (b)
*r1 = 5; *r2
\end{lstlisting}
Both fragments take two mutable references to \texttt{x}. Fragment
(a) writes through one of them and reads through the other; fragment
(b) does the same with the roles swapped. A static borrow check sees
the same set of borrows in scope at the same program points and has
to reject both. The runtime sees the actual sequence of operations
and can be more precise: the read after the write is safe in either
fragment, because the read does not modify state and the earlier
write through the other handle has not been challenged by a later
write. A small variant---\texttt{let r1 = r2;}---makes both handles
copies of one reference, and the question becomes uninteresting.

\paragraph{Two mutable arguments must not overlap.}
A function takes two mutable references. The caller must make sure
they refer to disjoint pieces of state, since otherwise the callee
cannot reason about either argument locally. Passing
\texttt{\&mut x.f} and \texttt{\&mut x.g} is fine---they touch
different fields. Passing \texttt{\&mut x} and \texttt{\&mut x.g} is
not---the second is reachable through the first. The runtime check is
substructural: it has to ask, for any two reference arguments, whether
one is reachable from the other through field projections,
indexing, and so on, along the actual paths the program took.

\medskip

The three examples look different on the surface: variant mutation,
mutable aliasing, and an overlap at a call boundary. But in every
case, one operation invalidates the path through which another
reference would read or write state. That observation is what
motivates the model. The next section describes the data structure
that records those paths; Sect.~\ref{sec:semantics} describes the
single rule that decides which references the next destructive event
invalidates.

\section{Access Path Trees}\label{sec:apt}

An access path describes how the program reached a particular
substructure of Move's state: it starts at some named entry
point---a local of the current frame, a globally-stored resource, or
a reference parameter passed in by the caller---and follows a
sequence of field projections, variant projections, and indexings.
Every borrow in the program produces a reference that can be
described by its access path. The runtime records these paths
explicitly: it groups them into a forest of trees, one per entry
point, and attaches each live reference to the node at the end of
its path.

A reference is \emph{valid} when the path it represents still leads
to a reachable, type-compatible substructure of the current state.
The runtime approximates this by marking references invalid when
destructive events happen along their paths; the rules in
Sect.~\ref{sec:semantics} spell out which events invalidate which
references.

\paragraph{Derivation.}
A reference $\refn{2}$ is \emph{derived from} $\refn{1}$ when its
access path is built by extending the path of $\refn{1}$ through one
of the language's borrowing operations: field projection, variant
projection, indexing, freezing, or being passed across a function
call as a reference parameter. In the access path tree this comes
out simply: $\refn{2}$ is derived from $\refn{1}$ exactly when its
node lies in the subtree below $\refn{1}$'s. The full derivation
order is the reflexive-transitive closure.

\paragraph{Roots and edges.}
Each tree's root is one of the named entry points mentioned above.
There are three kinds:
\begin{itemize}
\item $\local{i}$, the $i$-th local of the current frame;
\item $\glob{\tau}$, the global resource of type $\tau$ (we abstract
addresses);
\item $\param{i}$, the abstract value behind the $i$-th reference
parameter passed in by the caller.
\end{itemize}
Edges carry labels: a field offset for record fields, a variant
field for variant projections, and a single distinguished label $0$
for vector elements. The vector label is an over-approximation: rather
than track each index, we collapse all of them. This costs precision
on element-wise borrows but keeps the trees small.

Trees are built lazily. The runtime never materializes a node until
the program actually borrows through it.

\paragraph{State.}
A runtime state is a tuple $\sigma = (T, R, S, F)$.
\begin{itemize}
\item $T : \mathrm{Root} \rightharpoonup \mathrm{APT}$ maps each root
to its tree. The map is partial: roots that have not been borrowed
have no tree.
\item $R$ assigns each live reference identifier a triple $(\mu, \pi,
n)$. Here $\mu \in \{\mut, \imm\}$ is the mutability, $\pi \in
\{\Live, \Pois\}$ is the validity flag (we call $\Pois$ \emph{poisoned}
for short), and $n = \node{\mathrm{root}}{\mathrm{id}}$ is the
reference's \emph{qualified node}---the tree node it is attached to.
\item $S$ is a shadow stack mirroring the machine's operand stack.
Each entry is either a non-reference token or a reference identifier.
\item $F$ is the per-frame stack. Each frame carries its own $T$ and
$R$, a vector of shadow locals, and a map from this frame's
reference-parameter indices back to qualified nodes in the caller,
which is what return transport uses.
\end{itemize}

\paragraph{Aliasing as a tree relation.}
Two references alias when one's node is an ancestor or descendant of
the other's. The runtime decides this by walking parent pointers, in
time linear in the depth of the tree. Each node also holds the set of
references attached to it, so an invalidation walk only has to visit
the affected subtree---it does not have to scan the whole reference
table.

\section{Invalidation Rules}\label{sec:semantics}

We give the rules as a small-step relation $\sigma \to \sigma'$
indexed by the bytecode instruction being executed. Rules for
operations that do not touch the shadow state---constants,
arithmetic, stack reshuffles---are omitted.

Every non-trivial rule has the same shape. An event destructively
changes the substructure reachable through some node $n$, and the
runtime walks the tree from $n$ to poison every reference attached to
nodes related to $n$. Different events pick different sets of
related nodes (self, ancestors, descendants, or unions) and different
mutability filters. We write $\mathrm{invalidate}(N, \phi)$ for the
operation that poisons every reference attached to some node in
$N$, filtered by $\phi \in \{\mut, \imm, \mathsf{all}\}$. A poisoned
reference cannot be borrowed from, read, written, frozen, or passed
as a call argument; any such attempt is a runtime safety violation.

\paragraph{Borrow, read, freeze.}
Borrowing a local $i$ creates a fresh reference attached to
$\local{i}$, materializing $T(\local{i})$ if it does not yet exist. A
field borrow with label $\ell$ from a reference $\refn{}$ at node $n$
consumes $\refn{}$ and creates a fresh reference at the child
$n.\ell$. \texttt{ReadRef} pops the top reference after a validity
check; the node is unchanged. \texttt{FreezeRef} on a live $(\mut, n)$
reference consumes it and produces a fresh $(\imm, n)$ reference at
the same node---an immutable re-view of the same path. None of these
events invalidates anything.

\paragraph{Destructive write.}
\texttt{WriteRef} through a mutable reference at $n$ overwrites the
value at $n$, which affects both its ancestors (whose contents have
changed) and its descendants (which may have been replaced). The rule
is:
\[
\mathrm{invalidate}\bigl(\{n\} \cup \Desc(n) \cup \Anc(n),\ \imm\bigr),\qquad
\mathrm{invalidate}\bigl(\Desc(n),\ \mut\bigr).
\]
The first line poisons every immutable view of $n$, of any ancestor,
and of any descendant. The second poisons every mutable view of a
descendant. Mutable views of $n$ itself are deliberately spared: a
program that consumes one mutable view and immediately reborrows from
$n$ stays accepted, which matches static borrow checking.

\paragraph{Move, overwrite, variant mutation.}
Moving or storing over local $i$ replaces the value at $\local{i}$
outright, so we invalidate every reference reaching through it,
regardless of mutability:
$\mathrm{invalidate}(\{r_i\} \cup \Desc(r_i),\ \mathsf{all})$.
Variant mutation is the same rule applied at the discriminator:
references derived from the prior variant's projection sit in
descendant nodes, and the change to the discriminator poisons them.
The enum example from Sect.~\ref{sec:intro} falls out of this case.

\paragraph{Calls.}
A call $f(\overline{a})$ with reference arguments
$a_{i_1}, \ldots, a_{i_k}$, each at node $n_j$ with mutability
$\mu_j$, proceeds in four steps.
\begin{enumerate}
\item Check each $\refn{j}$ is live. Then lock the subtree at $n_j$:
exclusively if $\mu_j = \mut$, shared otherwise. An exclusive lock
conflicts with any other lock on a related node; a shared lock
conflicts only with an exclusive one. Any conflict is a violation.
This is what catches overlapping mutable arguments from
Sect.~\ref{sec:examples}.
\item For each mutable argument, apply the destructive-write rule at
$n_j$. The callee may write through the reference, so the discipline
treats the call itself as a write.
\item Allocate a new frame. The callee's reference parameters become
fresh references attached to parameter roots $\param{j}$. The new
frame records, for each $j$, the qualified node $n_j$ in the
caller---this is what return transport (below) uses.
\item Release the locks on the caller's subtrees.
\end{enumerate}
Fig.~\ref{fig:lock} illustrates a conflicting pair.
\begin{figure}[h]
\centering
\begin{minipage}{0.92\textwidth}\small
Calling \texttt{foo(r1, r2)} with $r_1 = \&\mut x$, $r_2 = \&\mut x.g$:
\begin{verbatim}
  root                       r1 takes exclusive lock on subtree(x).
   +-- x     <-- locked      r2 takes exclusive lock on subtree(x.g),
        +-- f                which is a descendant of x: lock conflict
        +-- g  <-- locked    ==> runtime safety violation.
\end{verbatim}
\end{minipage}
\caption{Call-boundary subtree locking.}\label{fig:lock}
\end{figure}

\paragraph{Returns.}
On return, every reference on the shadow stack must trace back to a
reference parameter of the current frame. Walking the parent pointers
of its qualified node should land on some $\param{j}$. A reference
rooted at a local or a global would dangle past the frame's lifetime
and is rejected. Each surviving reference is then transported into
the caller: starting from the caller's qualified node for parameter
$j$, the runtime walks the same labels and attaches a fresh reference
at the resulting node. Returning two mutable references that share a
subtree is caught by the same exclusive-lock pass used at calls.

\paragraph{Opaque effects.}
Some operations have bodies the runtime cannot see: native functions
implemented outside the bytecode language, closure captures, and
dynamic-dispatch targets. We handle them through a small interface:
each opaque operation declares, per returned reference, which input
reference argument it derives from. At a call site, the runtime
conjures a fresh child of that argument's node and attaches the
returned reference to it; ordinary call-boundary checks still apply.
The interface is part of the trusted base---an incorrect declaration
silently weakens the guarantees. We return to this in
Sect.~\ref{sec:discussion}.

\section{Precision against the Trace}\label{sec:precision}

The rules in Sect.~\ref{sec:semantics} are checked against the
program's actual execution, not against a join over the program's
possible executions. This is the discipline's main precision
advantage over the corresponding static check, and the reason we
expect a runtime model in this setting to be useful at all.

\paragraph{What the runtime knows that static analysis does not.}
The simplest illustration is a program with a conditional move:
\begin{lstlisting}
let x = ...;
let r = &x.f;
if rare_condition() { move x }   // invalidates subtree(x) only here
*r                               // safe whenever rare_condition() is false
\end{lstlisting}
A sound static analysis sees the move on one branch and has to
conclude that \texttt{r} may be invalid by the time \texttt{*r} runs;
it joins the two branches and reports the union. The discipline
sees something more specific. By the time control reaches \texttt{*r},
\texttt{rare\_condition()} has either evaluated to true or to false,
and the runtime knows which. On the branch where the move did not
happen, \texttt{r} is still attached to its node, no event has
poisoned it, and the read succeeds. The discipline accepts a
strictly larger set of executions than the static check, by
construction.

\paragraph{Relation to static checking.}
We are not trying to port a static borrow checker to runtime. The
rules in Sect.~\ref{sec:semantics} were calibrated so that each rule
poisons a reference only on cases that the corresponding static
analysis would also reject when reached on every branch; this is the
sense in which the discipline is meant to be at least as permissive
as its static counterpart along a concrete trace. We do not prove
that here. The natural proof obligation, for each static rule, is to
identify a runtime invariant strong enough to rule out the same
unsafe events. This is the most important open piece of the work,
discussed further in Sect.~\ref{sec:discussion}.

\paragraph{Where the cost is worth it.}
Per-instruction shadow-state updates are not free. The discipline
pays its way in contexts where static checking is unavailable or
untrusted---bytecode that bypasses the verifier (fuzzing harnesses,
mutated bytecode, dynamic-dispatch boundaries), and effects whose
body is opaque to the bytecode language (native code). On
verifier-protected execution paths, the discipline can be turned off
entirely, or shifted into a trace-and-replay mode that records the
relevant events during execution and checks them offline.

\section{Discussion and Future Work}\label{sec:discussion}

\paragraph{Formal soundness.}
The most important open piece is a small-step operational semantics
with a soundness theorem: if a program never uses a poisoned
reference, never has two overlapping mutable arguments at a call,
and never returns a reference that does not trace back to a
parameter, then its reference operations preserve type and memory
safety. The rules in Sect.~\ref{sec:semantics} are designed to make
this provable.

\paragraph{Opaque effects.}
The interface for native functions, closure captures, and
dynamic-dispatch targets is small but part of the trusted base. An
opaque operation declares which of its reference arguments each
returned reference is derived from; the runtime treats the operation
as a black-box step that produces a fresh child node under that
argument. This is enough to keep the discipline working at the
boundary, but correctness then depends on the declared model. A
natural follow-up is a unified treatment of all three cases together
with the soundness condition an opaque model must satisfy.

\paragraph{Deployment modes.}
Because invalidation is a one-way bit, the discipline can be deployed
in several modes. The most aggressive is \emph{always-on}, useful as
defense in depth on bytecode whose verifier outcome is not trusted.
A natural relaxation is \emph{trace-and-replay}: record the relevant
events during execution and run the check on the recorded trace
later, in parallel with other work. Production deployments will
likely mix the two depending on the trust they place in the static
verifier on each code path.

\paragraph{Limitations.}
Trusted models can silently break the guarantees if they misdescribe
an opaque effect. Vector indices collapse to a single abstract label,
so the discipline cannot distinguish two references into different
elements of the same vector. And when a violation is detected, the
runtime reports the event but does not yet emit a derivation trace
explaining which earlier operation poisoned the reference; such a
trace would help in debugging.

\section{Related Work}\label{sec:related}

\paragraph{Rust-style static borrow checking.}
The closest static relative is Rust's borrow checker, whose
metatheory was given by RustBelt~\cite{rustbelt}. Both target
stack-allocated state, with no first-class notion of globally-stored
mutable resources, and the only dynamic borrow check in mainstream
Rust is \texttt{RefCell}'s local check on a single cell. Move's
reference-safety rules, in contrast, span locals, globals, and
mutable cross-frame arguments; our discipline enforces them at
execution time.

\paragraph{Per-location dynamic alias models.}
The closest dynamic neighbors are Stacked Borrows~\cite{stackedborrows}
and Tree Borrows~\cite{treeborrows}. Stacked and Tree Borrows model
aliasing at memory locations: each location carries a stack or tree
of borrow tags, and the Miri interpreter checks reads and writes
against this metadata to detect undefined behavior in Rust. Our
model tracks references by their access paths rather than by the
addresses they happen to land on, which matters because Move's
reference rules include events the per-location view does not
naturally express---in particular variant mutation, where a write
to a discriminator invalidates references derived from the prior
variant's projection. We also treat call boundaries as a first-class
construct, with subtree locking and parameter-root transport doing
the work that the static borrow checker does at compile time in
Rust.

\paragraph{Capabilities, separation, runtime verification.}
Capability systems~\cite{capabilities} read naturally onto our model:
a node identity is the capability, and an invalidation is a
revocation. Access paths play the role of a runtime footprint in the
sense of separation logic~\cite{seplogic}, though we do not develop
the logical reading here. The RV literature on monitoring mutable
state~\cite{rvsurvey} provides the broader backdrop; we specialize
to one specific monitor, over the hierarchy of reference derivation.
Move ships with a static reference-safety checker performing abstract
interpretation over a borrow graph, which the runtime discipline is
designed to interoperate with.

\section{Conclusion}\label{sec:conclusion}

Move's static reference-safety rules can be enforced at execution
time using a small runtime discipline. The runtime records the
access paths through which each live reference was derived, and
applies a single invalidation rule when destructive events disturb
those paths; the rule covers writes, moves, freezes, variant
mutations, calls, returns, and opaque effects uniformly, and call
boundaries lock subtrees to keep mutable arguments disjoint. The
discipline catches a class of unsafe uses that lifetime-based static
analysis can only reject conservatively, and provides a uniform
handle on opaque effects at the boundary with native code. The most
important open work is a formal soundness proof against an
operational semantics; the model is designed to admit one.


\begin{credits}
\subsubsection{\discintname}
The authors have no competing interests relevant to this article.
\end{credits}

\bibliographystyle{splncs04}
\bibliography{biblio}

\end{document}

%%% Local Variables:
%%% mode: latex
%%% TeX-master: t
%%% End:
